\documentclass[12pt]{article}

\usepackage{url,verbatim, amsmath, amssymb, amsfonts, dsfont, color}
\textwidth 15true cm
\textheight 9true in
\oddsidemargin 0.30true in
\evensidemargin 0.30true in
\topmargin -0.3true in
\headsep 0.4true in

\def\sgn{{\rm sign}}
\def\drightarrow{{\stackrel{d}{\rightarrow}}}
\def\dotsim{\stackrel{.}{\sim}}
\def\goesd{\stackrel{d}{\rightarrow}}
\def\goesp{\stackrel{p}{\rightarrow}}
\def\goeslaw{\stackrel{{\cal L}}{\longrightarrow}}
%
\def\vp{{\varphi}}
\def\E{{\rm E}}
\def\Pvalue{{$p$-value}}
\def\Pvalues{{$p$-values}}
\def\pr{{\rm pr}}
\def\Pr{{\rm pr}}
\def\logit{{\rm logit}}
\def\T{{ \mathrm{\scriptscriptstyle T} }}
\def\diag{{\rm diag}}  
\def\half{{\textstyle{1\over 2}}}
\def\ith{{$i{\rm th}$ }}
\def\var{{\rm var}}
\def\cl{{{\rm c}\ell}}
\def\checkx{{\checked\hspace{-6pt}\setminus}}
\def\grad{{\nabla}}
\def\bs#1{{\boldsymbol{#1}}}

\sloppy %permits line-breaking of urls

\def\blue{\color{blue}}
\def\red{\color{red}}
\def\sgn{{\rm sign}}
\def\drightarrow{{\stackrel{d}{\rightarrow}}}
\begin{document}

\noindent{\bf Questions Week 5} {\hfill STA 2212S 2026}\\

%{\bf Due January 14} \hfill{\it MS, Exercise 5.2}

\bigskip

\begin{enumerate}

\item Suppose $X$ is a random variable with density $f(\cdot)$ and continuous distribution function $F(\cdot)$ defined over the whole real line $\mathbb{R}$.  Define the $p$-value for testing $H_0:F(\cdot) = F_0(\cdot)$ as $$
p^{obs}(x^{obs}) = \text{Pr}_{H_0}(X \ge x^{obs}),$$ where $x^{obs}$ is the observed value of $X$. 
\begin{enumerate}
\item 	Show that the distribution of $p^{obs}(X)$ is $U(0,1)$ under $H_0$. This question is a slight re-stating of {\blue AoS, Theorem 10.14}.
\item Suppose we carry out $k$ significance tests, leading to $p$-values $p_1^{obs}, \dots, p_k^{obs}$, and we report the smallest of these $q^{obs} = \min(p_j^{obs})$, say. One way to model this is to define the test statistic  
$$
Q = \min\{P_1, \dots, P_k\}, 
$$
where small values of $Q$ are evidence against $H_0$. 

Use the result of part (a) to show that
$$
\Pr(Q \le  q^{obs}; H_0) \le \sum_{j=1}^k\Pr(P_j \le q^{obs}) = kq^{obs}.
$$   
\end{enumerate}
\textcolor{red}{
(a) By definition, $p^{obs}(x) = P_{H_0}(X \geq x) = 1 - F_0(x)$. Let $p^* \in [0, 1]$. Since the $p$-value is a probability, it must lie within $[0, 1]$, so the cdf of the $p$-value can be computed as
\begin{align*}
    P_{H_0}(p^{obs}(X) \leq p^*) &= P_{H_0}(1 - F_0(X) \leq p^*) \\
    &= P_{H_0}(X \geq F_0^{-1}(1 - p^*)) \\
    &= 1 - F_0(F_0^{-1}(1 - p^*))\\
    &= p^*
\end{align*}
thus the cdf of $p^{obs}$ is the cdf of a $U(0, 1$) random variable under $H_0$. So, $p^{obs} \sim_{H_0} U(0, 1)$.
}

\noindent \textcolor{red}{
(b) By finite subadditivity, 
\begin{align*}
    P_{H_0}(Q \leq q^{obs}) &= P_{H_0}(\min\{P_1, \ldots, P_k\} \leq q^{obs})\\
    &= P_{H_0} \left(\bigcup_{i=1}^k \{P_i \leq q^{obs}\}\right)\\
    &\leq \sum_{i=1}^k P_{H_0} (P_i \leq q^{obs})\\
    &= \sum_{i=1}^k q^{obs}\\
    &= kq^{obs}
\end{align*}
}
\item Suppose that $(X_{1i}, X_{2i}), i = 1, \dots, n$ follow the bivariate normal distribution with expected value $(\theta_1, \theta_2)$ and identity covariance matrix: the joint distribution of the sample $(x_1, x_2) = (x_{11}, \dots , x_{1n}, x_{21}, \dots, x_{2n})$ is then
$$
f( x_1,  x_2;\theta) = \left(\frac{1}{2\pi}\right)^n\exp\left[-\frac{1}{2}\left\{\sum(x_{1i}-\theta_1)^2 + \sum(x_{2i}-\theta_2)^2\right\}\right]
$$
The parameter space is restricted to $\theta_1 \ge 0, \theta_2 \ge 0$.
\begin{enumerate}

\item Show that the maximum likelihood estimators of $\theta_1$ and $\theta_2$ are given by
$$
\hat\theta_1 = \max(\bar x_1,0), \quad \hat\theta_2 = \max(\bar x_2, 0).
$$
\item  Derive the form of the log-likelihood ratio statistic $w = 2\{\ell(\hat\theta) - \ell(\theta_0)\}$ for testing $H_0: (\theta_1, \theta_2) = (\theta_{01},\theta_{02}) = (0,0)$. %{\blue It has a different expression for $(\bar x_1, \bar x_2)$ in each of the 4 quadrants of the plane.}
\item   Show that 
the distribution of $w$ under $H_0$ is 
$$
\frac{1}{4}\delta_{\{0\}} + \frac{1}{2}\chi_1^2 + \frac{1}{4}\chi_2^2,
$$
where $\delta_{\{0\}}$ is a point mass at $0$.
\end{enumerate}
\end{enumerate}
\textcolor{red}{
(a) The log-likelihood is 
\begin{equation*}
    \ell(\theta_1, \theta_2 ; x_1, x_2) = -\frac{1}{2}\left[\sum_{i=1}^n \left((x_{1i} - \theta_1)^2 + (x_{2i} - \theta_2)^2 \right)\right] + C
\end{equation*}
hence 
\begin{align*}
    \frac{\partial \ell}{\partial \theta_1} &= \sum_{i=1}^n (\theta_1 - x_{1i}) = n\theta_1 - \sum_{i=1}^n x_{1i}\\
    \frac{\partial \ell}{\partial \theta_2} &= \sum_{i=1}^n (\theta_2 - x_{2i}) = n\theta_2 - \sum_{i=1}^n x_{2i}
\end{align*}
Setting both to 0 and solving for $\theta_1$ and $\theta_2$ over the parameter space yields MLEs 
\begin{equation*}
    \hat\theta_1 = \max(\bar x_1, 0) \qquad \hat\theta_2 = \max(\bar x_2, 0)
\end{equation*}
since $\ell$ is decreasing in both $\theta_1$ and $\theta_2$. 
}

\noindent \textcolor{red}{
(b) From part(a), 
\begin{equation*}
    \ell(\hat\theta) = -\frac{1}{2}\left[\sum_{i=1}^n \left((x_{1i} - \hat\theta_1)^2 + (x_{2i} - \hat\theta_2)^2 \right)\right] + C
\end{equation*}
and 
\begin{equation*}
    \ell(\theta_0) =  -\frac{1}{2}\left[\sum_{i=1}^n \left(x_{1i}^2 + x_{2i}^2 \right)\right] + C
\end{equation*}
Then 
\begin{align*}
    (x_{1i} - \hat\theta_1)^2 + (x_{2i} - \hat\theta_2)^2 - x_{1i}^2 - x_{2i}^2 &= -2x_{1i}\hat\theta_1 - 2x_{2i}\hat\theta_2 + \hat\theta_1^2 + \hat\theta_2^2\\
    &= \hat\theta_1(\hat\theta_1 - 2x_{1i}) + \hat\theta_2(\hat\theta_2 - 2x_{2i})
\end{align*}
This implies 
\begin{align*}
    w &= 2\left(-\frac{1}{2}\sum_{i=1}^n (\hat\theta_1(\hat\theta_1 - 2x_{1i}) + \hat\theta_2(\hat\theta_2 - 2x_{2i}))\right)\\
    &= \sum_{i=1}^n \left(\hat\theta_1(2x_{1i} - \hat\theta_1) + \hat\theta_2(2x_{2i} - \hat\theta_2)\right) \\
    &= \hat\theta_1 \left(2\sum_{i=1}^n x_{1i} - n\hat\theta_1\right) + \hat\theta_2\left(2\sum_{i=1}^n x_{2i} - n\hat\theta_2\right)
\end{align*}
If $\bar x_1 > 0$, then 
\begin{equation*}
    \hat\theta_1 \left(2\sum_{i=1}^n x_{1i} - n\hat\theta_1\right) = n\bar x_1^2
\end{equation*}
and similarly if $\bar x_2 > 0$, then 
\begin{equation*}
    \hat\theta_2 \left(2\sum_{i=1}^n x_{2i} - n\hat\theta_2\right) = n\bar x_2^2
\end{equation*}
However if $\bar x_1 \leq 0$, then 
\begin{equation*}
    \hat\theta_1 \left(2\sum_{i=1}^n x_{1i} - n\hat\theta_1\right) = 0
\end{equation*}
and similarly if $\bar x_2 \leq 0$, then 
\begin{equation*}
    \hat\theta_2 \left(2\sum_{i=1}^n x_{2i} - n\hat\theta_2\right) = 0
\end{equation*}
thus 
\begin{equation*}
    w =n \left(\max(\bar x_1, 0)^2 + \max(\bar x_2 , 0)^2\right)
\end{equation*}
}

\noindent \textcolor{red}{
(c) We consider cases $\bar x_1 \leq 0, \bar x_2 > 0$, $\bar x_1 > 0, \bar x_2 \leq 0$, $\bar x_1, \bar x_2 \leq 0$, and $\bar x_1, \bar x_2 > 0$. Then 
\begin{equation*}
    w = \begin{cases}
        0 & \bar x_1 \leq 0, \bar x_2\leq 0\\
        n\bar x_1^2 & \bar x_1 > 0, \bar x_2 \leq 0\\
        n\bar x_2^2 & \bar x_1 \leq 0, \bar x_2 > 0\\
        n(\bar x_1^2 + \bar x_2^2) & \bar x_1 > 0, \bar x_2 > 0
    \end{cases}
\end{equation*}
Under the null, we have $(X_{1i}, X_{2i}) \sim \mathcal N_2(0, I_2)$, hence 
\begin{equation*}
    \bar X_1 \sim_{H_0} \mathcal N\left(0, \frac 1n\right) \qquad \bar X_2 \sim_{H_0} \mathcal N\left( 0, \frac 1n \right)
\end{equation*}
Then 
\begin{align*}
    P(w = 0) &= P(\bar X_1 \leq 0, \bar X_2 \leq 0) = \frac{1}{4}\\
    P(w = n\bar X_1^2) &= P(\bar X_1 > 0, \bar X_2 \leq 0) = \frac 14\\
    P(w = n\bar X_2^2) &= P(\bar X_1 \leq 0, \bar X_2 > 0) = \frac 14\\
    P(w = n\bar X_1^2 + n\bar X_2^2) &= P(\bar X_1 > 0, \bar X_2 > 0) = \frac 14
\end{align*}
Then the distribution of $w$ under $H_0$ is 
\begin{equation*}
    \frac{1}{4}\delta_{\{0\}} + \frac 12 \chi_1^2 + \frac 14 \chi_2^2
\end{equation*}
}
\end{document}
